\section{Maximum Bracket Size}

Dalam sistem eliminasi, ukuran bracket harus mengikuti bilangan pangkat dua. 
Hal ini dilakukan agar setiap ronde dapat membagi jumlah peserta secara seimbang hingga tersisa satu pemenang.

Ukuran bracket yang umum digunakan adalah:

\[
2^n \quad \text{dengan } n \in \mathbb{N}
\]

Sehingga ukuran bracket yang valid antara lain:

\[
2, 4, 8, 16, 32, 64, 128, 256
\]

\subsection{Penentuan Maximum Bracket}

Maximum bracket size ditentukan sebagai bilangan pangkat dua terkecil yang lebih besar atau sama dengan jumlah peserta.

Secara matematis:

\[
\text{Bracket Size} = 2^{\lceil \log_2(N) \rceil}
\]

dengan:

\begin{itemize}
\item $N$ = jumlah peserta
\item $\lceil \cdot \rceil$ = fungsi pembulatan ke atas (ceiling)
\end{itemize}

\subsection{Contoh Perhitungan}

Jika jumlah peserta adalah:

\begin{itemize}
\item $N = 60$
\end{itemize}

maka:

\[
\log_2(60) \approx 5.91
\]

\[
\lceil 5.91 \rceil = 6
\]

Sehingga:

\[
\text{Bracket Size} = 2^6 = 64
\]

Artinya bracket eliminasi yang digunakan adalah \textbf{64 slot}.

\subsection{Hubungan dengan Bye}

Jika jumlah peserta lebih kecil dari ukuran bracket, maka slot kosong akan menjadi \textbf{bye}.

Rumus jumlah bye:

\[
\text{Bye} = \text{Bracket Size} - N
\]

Contoh:

\[
64 - 60 = 4
\]

Sehingga terdapat \textbf{4 bye} pada ronde pertama.

\subsection{Contoh Tabel Penentuan Bracket}

\begin{center}
\begin{tabular}{|c|c|c|}
\hline
Jumlah Peserta & Bracket Size & Bye \\
\hline
14 & 16 & 2 \\
27 & 32 & 5 \\
60 & 64 & 4 \\
70 & 128 & 58 \\
\hline
\end{tabular}
\end{center}

\subsection{Kesimpulan}

\begin{itemize}
\item Ukuran bracket eliminasi selalu mengikuti bilangan pangkat dua.
\item Maximum bracket size dipilih sebagai pangkat dua terkecil yang lebih besar atau sama dengan jumlah peserta.
\item Jika jumlah peserta tidak memenuhi ukuran bracket, maka slot kosong diisi dengan bye.
\end{itemize}